\newcommand{\tipoRequerimento}{Taxa administrativa da CIP}

\section{DOS FATOS E DOS DIREITOS}

\begin{enumerate}[resume]
    \item A concessionária indeferiu a reclamação referente a taxa administrativa da CIP do município, sob o fundamento de que as cobranças questionadas ocorreram em conformidade com o contexto normativo, contratual e judicial vigente à época.

    \item Nesse contexto, a controvérsia entre o Município e a Concessionária não se limita à paralisação da cobrança da taxa de administração da CIP/COSIP, obrigação já cumprida após o Sexto Termo Aditivo, mas centra-se na restituição dos valores arrecadados no passado. A concessionária recusa-se a devolver o montante, alegando que as cobranças tinham respaldo judicial pretérito e que o novo regime não possui efeitos retroativos.

    \item Contudo, a demanda do Município não perdeu seu objeto. O argumento principal é que a licitude pretérita da cobrança (baseada em decisões precárias) não gera direito adquirido à retenção definitiva dos valores. Ao aderir voluntariamente ao \termoCIP{} para garantir a prorrogação da concessão, a Concessionária aceitou as novas diretrizes da ANEEL sem reservas, o que incluiu a obrigatoriedade de renunciar a ações judiciais (próprias ou via associações como a ABRADEE).

    \item A tentativa da concessionária de cumprir o acordo apenas parcialmente, desistindo das ações, mas retendo os lucros passados decorrentes delas, esvazia a finalidade da cláusula de renúncia. Tal postura caracteriza resistência à autoridade regulatória e viola frontalmente a boa-fé objetiva, a segurança jurídica e a vedação ao comportamento contraditório (\textit{venire contra factum proprium}).

    \item Em suma, a adesão ao novo contrato impede a preservação seletiva de vantagens econômicas fundadas em teses jurídicas que a própria empresa se obrigou a abandonar, tornando devida a recomposição patrimonial ao Município.

    \item Ressalta-se, ainda, que a Ouvidoria da ANEEL deve observar e cumprir integralmente o entendimento consignado no Parecer nº 103/2026 (em anexo), emitido pela própria Procuradoria da Agência, de modo que a análise da presente demanda seja realizada em conformidade com as orientações ali estabelecidas.

\end{enumerate}

\section{DOS PEDIDOS}

\begin{enumerate}[resume]
\item Ante o exposto, requer-se à Concessionária:
    \begin{enumerate}
        \item Proceda à devolução dos valores indevidamente cobrados a maior a título de taxa de administração pela arrecadação da CIP, referentes ao período de 13/10/2020 até a data da efetiva e definitiva cessação da cobrança;

        \item O cálculo dos valores referidos no item anterior deverá observar, rigorosamente, os seguintes parâmetros:
        
        \begin{enumerate}
            \item no período compreendido entre 13/10/2020 e 29/04/2021, a cobrança da taxa de administração deverá ser limitada a 1\% (um por cento) do montante efetivamente arrecadado, sempre que o percentual contratualmente pactuado superar tal patamar; e 
            \item A partir de 29/04/2021, o valor devido a esse título deve ser igual a zero, isto é, a cobrança da taxa de administração da CIP/COSIP deveria ter sido integralmente cessada naquela data;
        \end{enumerate}

        \item Que todas as comunicações, intimações e notificações relativas à presente reclamação sejam encaminhadas exclusivamente ao representante legal do Município que a subscreve, preferencialmente por meio do endereço eletrônico \emailEmpresa{}, em observância aos princípios do contraditório, ampla defesa, eficiência e celeridade, bem como ao art. 9º da REN ANEEL nº 1.000/2021.

    \end{enumerate}
\end{enumerate}